\section{Normalizer extension and Frobenius generation}\label{sec:generation}

Let $F$ be a number field and let
\[
  \alpha_0:G_F\longrightarrow P
\]
be a continuous open homomorphism. Our eventual aim is to identify this
map with ordinary restriction, up to conjugation in $P$. We first enlarge
an open part of its domain in a canonical way. For the enlarged map we
will prove surjectivity and equality of cyclotomic logarithms at an
infinite set of auxiliary primes. The main step is to show that images
of local Frobenius elements generate prescribed cyclic subgroups in
suitable finite quotients of $P$.

Choose an open normal subgroup $U\trianglelefteq\mathcal G$ contained in
$G_F$, and put
\[
 A=\alpha_0(U),\qquad N=\ker(\alpha_0|_U),\qquad
 T=N_{\mathcal G}(N).
\]
The subgroup $A$ is open in $P$. Since $U\subset T$, the normalizer $T$
is open and closed in $\mathcal G$; write $T=G_k$.

\begin{lemma}[Extension to the normalizer]\label{lem:normalizer-extension}
There is a unique continuous open homomorphism
\[
  \beta:T\longrightarrow P
\]
whose restriction to $U$ is $\alpha_0|_U$.
\end{lemma}

\begin{proof}
Conjugation by $t\in T$ induces a continuous automorphism of $U/N=A$.
By Theorem~\ref{thm:uchida-isomorphism}, applied to
$\Q^s/(\Q^s)^A$, this automorphism is conjugation by a unique element
$\beta(t)$ of $P$. Indeed, the implementing field automorphism of
$\Q^s$ fixes $\Q$. Uniqueness shows both that $\beta$ is a homomorphism
and that $\beta(u)=\alpha_0(u)$ for $u\in U$. It also shows that any
extension of $\alpha_0|_U$ to $T$ must equal $\beta$: its values must
induce these same conjugations on $A$.

Continuity on the open subgroup $U$ implies continuity on $T$.
Moreover, $\beta(T)$ contains $A$ and is therefore open in $P$.
A continuous homomorphism of profinite groups is open onto its image,
so $\beta$ is open.
\end{proof}

We keep $U,A,N,T,k$, and $\beta$ fixed for the rest of the argument.
List the distinct conjugates of $N$ other than $N$ itself as
\[
 N_i=t_iNt_i^{-1}\quad(1\leq i\leq r),\qquad t_i\notin T,
 \qquad H_i=\beta(N_i).
\]
There are finitely many because $U$ normalizes $N$. Normality of $U$
in $\mathcal G$ implies that each $N_i$ lies in $U$ and is normal in
$U$, so $H_i$ is normal in $A$.
Each $H_i$ is nontrivial. Otherwise $t_iNt_i^{-1}\subset N$, and
iterating this inclusion would give a descending chain returning to
$N$: some power of $t_i$ lies in $U$ and hence normalizes $N$.
All inclusions would be equalities, contrary to $t_i\notin T$.
We allow $r=0$, with sums indexed by $i$ then equal to zero.

\subsection{A finite quotient with enough characters}

The next lemma constructs characters that distinguish a proper
subgroup of a cyclic group while avoiding the invariant spaces attached
to the $H_i$. The latter condition will prevent cancellation by
conjugates of $N$ when we average a pulled-back character on the source.
For a finite subgroup $J$ acting on a characteristic-zero vector space,
we use the averaging operator $E_J=|J|^{-1}\sum_{g\in J}g$.

\begin{lemma}[Quotient amplification]\label{lem:quotient-amplification}
Fix an odd prime $\ell$ and a finite quotient $R_0$ of $P$. There is
a finite Galois extension $L/\Q$ in $\Q^s$, totally imaginary, such that
the quotient $R=\Gal(L/\Q)$ dominates $R_0$ and has the following
property. For every $\gamma\in R$, with $C=\langle\gamma\rangle$, and
every proper subgroup $D<C$, there is $x\in\X{\ell}{L}$ satisfying
\begin{equation}\label{eq:separating-character}
 x\in\X{\ell}{L}^{D},\qquad E_Cx=0,\qquad
 x\notin\sum_{i=1}^r\X{\ell}{L}^{H_i}.
\end{equation}
Here each $H_i$ acts through its image in $R$.
\end{lemma}

\begin{proof}
Choose a quotient $R_1=\Gal(L_1/\Q)$ dominating $R_0$ such that
every $H_i$ has nontrivial image in $R_1$, the kernel of $P\to R_1$
is contained in $A$, and $L_1$ contains $\Q(i)$. These are finitely
many compatible conditions on an open normal subgroup of $P$.
Put $n=|R_1|$.

We will enlarge $R_1$ while keeping element orders bounded by $2n$.
At the same time, the images of the $H_i$ and the conjugacy class of
complex conjugation will grow. Growth of the $H_i$ makes the sum
of their invariant spaces small. Growth of the conjugacy class
makes the signed representation supply enough $D$-invariant
vectors killed by $E_C$ to escape that sum.

We use the regular-module auxiliary extension construction appearing
in Uchida's proof~\cite[pp.~360--362]{uchida1979}, with quantitative
bounds on its invariant spaces. More precisely, we construct a
further quotient with kernel consisting of $m$ copies of the regular
$\mathbb F_2[R_1]$-module, where $m$ will be chosen below.
Choose $m$ distinct rational primes that split completely in $L_1$.
For each of them, select one prime of $L_1$ above it. Approximation
gives elements $a_1,\ldots,a_m\in L_1^\times$ such that $a_j$ has
valuation one at its selected prime and valuation zero at every other
prime in these $m$ sets. The $mn$ elements
\[
  g(a_j),\qquad g\in R_1,\quad 1\leq j\leq m,
\]
have independent square classes: their valuations at the selected
sets of primes form a permutation of the identity matrix modulo two.
Adjoin their square roots and call the resulting field $L$.
Kummer theory gives
\begin{equation}\label{eq:regular-kummer-kernel}
  1\longrightarrow V\longrightarrow R=\Gal(L/\Q)
   \longrightarrow R_1\longrightarrow1,
  \qquad V\simeq\mathbb F_2[R_1]^m
\end{equation}
as $R_1$-modules. The dual module supplied by Kummer theory is again
the regular permutation module. The field $L$ is Galois over $\Q$
because all conjugates of the $a_j$ were included; its group is
solvable, so $L\subset\Q^s$. No splitting of
\eqref{eq:regular-kummer-kernel} is required.

Every element of $R$ has order at most $2n$, since its $n$th power
lies in the exponent-two group $V$. Write $A_R$ and $H_{i,R}$ for
the images of $A$ and $H_i$ in $R$. The choice of $R_1$ gives
$V\subset A_R$. Choose $h_i\in H_{i,R}$ with nonidentity image
$g_i\in R_1$, of order $d_i>1$. As $H_{i,R}$ is normal in $A_R$,
\[
 (g_i-1)V=[h_i,V]\subset H_{i,R}.
\]
On $m$ regular modules the rank of $g_i-1$ is
$mn(1-1/d_i)\geq mn/2$. Consequently
\begin{equation}\label{eq:large-subgroups}
  |H_{i,R}|\geq 2^{mn/2}.
\end{equation}
If $c\in R$ is a complex conjugation, its image in $R_1$ is a
nontrivial involution because $L_1$ is totally imaginary.
Conjugation by $V$ alone therefore gives at least $2^{mn/2}$ distinct
conjugates of $c$; thus
\begin{equation}\label{eq:large-conjugacy-class}
  |c^R|\geq 2^{mn/2}.
\end{equation}

Set $\eta=2^{-mn/2}$ and write $X=\X{\ell}{L}$. By
Proposition~\ref{prop:signed-representation}, $X$ embeds in the regular
$\Q_\ell[R]$-representation and contains
\[
 W=\Ind_{\langle c\rangle}^{R}(\operatorname{sign}).
\]
The regular-representation bound and \eqref{eq:large-subgroups} give
\begin{equation}\label{eq:forbidden-dimension}
 \dim\left(\sum_{i=1}^rX^{H_i}\right)
 \leq\sum_{i=1}^r\frac{|R|}{|H_{i,R}|}
 \leq r\eta|R|.
\end{equation}
The character of $W$ takes the value $|R|/2$ at the identity,
the value $-|R|/(2|c^R|)$ on $c^R$, and zero elsewhere.
Hence, for every subgroup $J\subset R$,
\begin{equation}\label{eq:signed-invariants}
 \left|\dim W^J-\frac{|R|}{2|J|}\right|
 \leq\frac{\eta|R|}{2}.
\end{equation}
For $D<C=\langle\gamma\rangle$, the map
$E_C:W^D\to W^C$ is surjective. Since $|D|\leq |C|/2$ and
$|C|\leq2n$, \eqref{eq:signed-invariants} implies
\begin{align*}
 \dim\ker(E_C|_{W^D})
 &=\dim W^D-\dim W^C\\
 &\geq\frac{|R|}{2}\left(\frac1{|D|}-\frac1{|C|}\right)
       -\eta|R|\\
 &\geq |R|\left(\frac1{4n}-\eta\right).
\end{align*}
Choose the integer $m\geq1$ so large that
\[
  (r+1)2^{-mn/2}<\frac1{4n}.
\]
This single choice works for all $\gamma$ and all proper $D<C$.
The last dimension is then greater than
\eqref{eq:forbidden-dimension}, so its kernel contains a vector
outside the indicated sum. This vector gives
\eqref{eq:separating-character}. If $C$ is trivial, there is no
proper subgroup to consider.
\end{proof}

\subsection{Generation by local Frobenius images}

We now use these characters to strengthen the local correspondence.
Containment of a source decomposition group in a target decomposition
group places its Frobenius image in a cyclic subgroup of a finite
quotient. The next proposition shows that, after enlarging the
quotient, infinitely many such images generate the entire cyclic
subgroup.

\begin{proposition}[Frobenius generation]\label{prop:frobenius-generation}
Fix an odd prime $\ell$, a finite quotient $R_0$ of $P$, and
$\gamma_0\in R_0$. There are a quotient
\[
 P\twoheadrightarrow R=\Gal(L/\Q)\twoheadrightarrow R_0
\]
and a lift $\gamma\in R$ of $\gamma_0$ with the following property.
There are infinitely many $\ell$-active primes $v$ of $k$, over
distinct rational primes $p\ne\ell$, for which compatible
prolongations can be chosen so that:
\begin{enumerate}
\item the corresponding target rational prime $q\ne\ell$ is unramified
in $L$, and the target prolongation restricts to a prime $w'$ of $L$
whose Frobenius in $R$ is $\gamma$;
\item the image in $R$ of an arithmetic Frobenius lift at $v$ generates
$\langle\gamma\rangle$.
\end{enumerate}
\end{proposition}

\begin{proof}
Choose $R$ by Lemma~\ref{lem:quotient-amplification} and choose any
lift $\gamma$ of $\gamma_0$. Put $C=\langle\gamma\rangle$.
Take an open normal subgroup $U'\trianglelefteq\mathcal G$ contained
in $U\cap\beta^{-1}(P_L)$; for example, take the core of this open
subgroup in $\mathcal G$. Let $F'/\Q$ be the finite Galois extension
with $G_{F'}=U'$. The map
\[
 \beta|_{U'}:G_{F'}\longrightarrow P_L
\]
is open.

Chebotarev's theorem supplies infinitely many primes $w'$ of $L$,
over distinct rational primes $q\ne\ell$ unramified in $L$, with
Frobenius exactly $\gamma$. Theorem~\ref{thm:local-correspondence}
applied to $\beta|_{U'}$ pairs all but finitely many of them with
$\ell$-active primes $v'$ of $F'$, away from $\ell$.
Uniqueness of the target prime for an active source prime implies
that infinitely many $v'$ occur. Since only finitely many primes of
$F'$ lie over a given rational prime, we may pass to a subset with
distinct source rational primes $p$, all unramified in $F'$.
Passing to another infinite subset, we may also assume that their
Frobenius elements in $\Gal(F'/\Q)$ are a fixed element $\delta$.

Fix a source prolongation above each $v'$, and use it also for the
prime $v=v'|_k$. The original correspondence at $v'$ can be made
with this prolongation: changing the source prolongation conjugates
the target group by an element of $\beta(U')\subset P_L$, which
preserves the finite prime $w'$. Activity passes from $v'$ to $v$,
since the relevant inertia subgroup for $U'$ is contained in that
for $T$. Apply Theorem~\ref{thm:local-correspondence} at $v$ using
this same source prolongation. If $w_1$ is the restriction to $L$
of its target prolongation, then
\[
 \beta(D_{v'})\subset\beta(D_v)\cap P_L\subset D_{w_1}.
\]
Uniqueness of the target prime for $v'$ at the level of $L$ gives
$w_1=w'$. Consequently the corresponding target rational prime
is $q$, and the image of $D_v$ in $R$ lies in $C$.
These changes of prolongations preserve $v'$ and $w'$, so they
preserve the previously selected finite Frobenius elements
$\delta$ and $\gamma$.

Put $H=T/U'\subset\mathcal G/U'$ and let $f$ be the residue degree
of $v$ over $p$. Figure~\ref{fig:generation-primes} records the
primes involved. Here $\delta=\Fr(v'/p)$ and
$\gamma=\Fr(w'/q)$ belong to the finite Galois groups over $\Q$.
The symbols $\Fr_{v'}$ and $\Fr_{w'}$ used below instead denote
local arithmetic Frobenius over $F'$ and $L$, respectively.

\begin{figure}[htbp]
\centering
\begin{tikzcd}[row sep=large,column sep=huge]
 v'\text{ of }F' \arrow[d] \arrow[r,dashed] &
 w'\text{ of }L \arrow[d] \\
 v\text{ of }k \arrow[d,"f"'] \arrow[r,dashed] &
 q\text{ of }\Q \\
 p\text{ of }\Q &
\end{tikzcd}
\caption{Vertical arrows restrict primes; dashed arrows are the
partial local correspondences. The residue degree of $v/p$ is $f$.
No equality between $p$ and $q$ is assumed at this stage.}
\label{fig:generation-primes}
\end{figure}

In the decomposition group
$\langle\delta\rangle\subset\Gal(F'/\Q)$ we have
\begin{equation}\label{eq:residue-degree-intersection}
 \langle\delta\rangle\cap H=\langle\delta^f\rangle,
 \qquad
 f=\min\{j\geq1:\delta^j\in H\}.
\end{equation}
Indeed, the residue degree is the orbit size of the distinguished
coset in the action of $\langle\delta\rangle$ on
$\Gal(F'/\Q)/H$. This description does not require $H$ to be
normal. In particular $f$ is constant on the chosen infinite set.
The homomorphism $\beta$ induces $H\to R$, since
$\beta(U')\subset P_L$; define
\[
 d=\beta(\delta^f)\in R.
\]
This notation is independent of the lift of $\delta^f$ to $T$.
A local lift of rational Frobenius to the $f$th power is a
Frobenius lift at $v$, so the preceding containment gives $d\in C$.

We claim that $D=\langle d\rangle$ equals $C$. Suppose that $D<C$,
and choose $x$ as in \eqref{eq:separating-character}. Pullback along
$\beta|_{U'}$ gives a linear map
\[
 j:\X{\ell}{L}\longrightarrow\X{\ell}{F'},\qquad y=j(x).
\]
The map $j$ is injective: its defining homomorphism has open image
in $P_L$, and an additive $\Q_\ell$-valued character vanishing on
an open subgroup vanishes everywhere. Frobenius evaluation at $w'$
is invariant under $\gamma$, because $\gamma$ fixes $w'$ and acts
trivially by conjugation on its unramified local quotient. Therefore
\[
 x(\Fr_{w'})=(E_Cx)(\Fr_{w'})=0.
\]
The character $x$ is unramified at $w'$, so this equality says that
it vanishes on the entire decomposition group there. It follows
that $y(\Fr_{v'})=0$ for every prime in our infinite set.
Proposition~\ref{prop:character-detection}(ii) now gives
\begin{equation}\label{eq:source-average-zero}
 E_{\langle\delta\rangle}y=0.
\end{equation}

To see why this contradicts the choice of $x$, write
$a=|\langle\delta\rangle|$ and $Y=j(\X{\ell}{L})$.
Pullback is $T$-equivariant, where the action on the target
character space is through $\beta$. In particular $y$ is fixed
by $N$. In the sum defining \eqref{eq:source-average-zero}, the
powers $\delta^j$ belonging to $H$ are exactly those for which
$f\mid j$, by \eqref{eq:residue-degree-intersection}. Each such
power acts as the identity on $y$, because $x$ is $D$-invariant.
If $f\nmid j$, choose a lift $t$ of $\delta^j$ to $\mathcal G$.
Then $t\notin T$, and $\delta^jy$ is fixed by $tNt^{-1}$.
Multiplying \eqref{eq:source-average-zero} by $a$ thus yields
\[
  0=\frac af\,y+\sum_{\substack{0\leq j<a\\f\nmid j}}\delta^jy,
  \qquad
  y\in\sum_{i=1}^r\X{\ell}{F'}^{N_i}.
\]

It remains to return this membership to the target character
space. The finite group $U/U'$ acts on $\X{\ell}{F'}$, and its
subspace $Y$ is stable. Semisimplicity supplies a $U$-equivariant
projection
\[
  \pi:\X{\ell}{F'}\longrightarrow Y.
\]
Every $N_i$ lies in $U$, so $\pi$ carries $N_i$-invariants into
$Y^{N_i}$. Equivariance and injectivity of $j$ give
\[
 Y^{N_i}=j\bigl(\X{\ell}{L}^{H_i}\bigr).
\]
Applying $\pi$ to the displayed membership and then applying
$j^{-1}$ contradicts the last condition in
\eqref{eq:separating-character}. Only $U$-equivariance of the
projection is used; no normality of $T$ in $\mathcal G$ is needed.
Thus $D=C$, proving the proposition.
\end{proof}

\subsection{Surjectivity and cyclotomic logarithms}

The generation statement now gives the two properties needed in the
next stage. The first removes any restriction on the target image.
For the second, a congruence condition on $\ell$ forces the source
primes supplied by the proposition to have residue degree one, so
the character-detection theorem applies.

\begin{theorem}\label{thm:normalized-map}
The extension $\beta:G_k\to P$ is surjective. Put
\[
 m_0=[k:\Q]!,\qquad
 \mathcal L=\{\ell:\ell\text{ is an odd prime and }
                    \ell\equiv1\pmod{m_0}\}.
\]
The set $\mathcal L$ is infinite, and for every $\ell\in\mathcal L$,
\begin{equation}\label{eq:cyclotomic-logarithms}
 \log\chi_\ell\circ\beta=\log\chi_{k,\ell}
 \quad\text{as additive characters on }G_k.
\end{equation}
Here $\chi_\ell$ is the $\ell$-adic cyclotomic character on $P$,
and $\chi_{k,\ell}$ is the ordinary cyclotomic character on $G_k$.
\end{theorem}

\begin{proof}
Let $R_0$ be any finite quotient of $P$ and let $\gamma_0\in R_0$.
Proposition~\ref{prop:frobenius-generation} produces an element
$d$ in the image of $\beta$ in a further quotient $R$ such that
$\langle d\rangle=\langle\gamma\rangle$, where $\gamma$ lifts
$\gamma_0$. Thus $\gamma$, and hence $\gamma_0$, belongs to the
image. Since $\beta(T)$ is compact, surjectivity onto every finite
quotient implies $\beta(T)=P$.

Dirichlet's theorem on primes in arithmetic progressions shows
that $\mathcal L$ is infinite. Fix $\ell\in\mathcal L$, take
\[
 R_0=\Gal(\Q(\mu_\ell)/\Q)\simeq\mathbb F_\ell^\times,
\]
and choose $\gamma_0$ to generate this cyclic group.
The proposition gives infinitely many active primes $v$ of $k$
over distinct rational primes $p\ne\ell$, for which the image of
a Frobenius lift $\sigma_v$ acts on $\mu_\ell$ by a generator.
Let $q$ be the corresponding target rational prime and let
$f=[\kappa(v):\mathbb F_p]$, so $Nv=p^f$. If
$a(\sigma_v)\in\widehat{\Z}$ is the target residue exponent,
Lemma~\ref{lem:norm-identity} gives
\begin{equation}\label{eq:generating-norm}
  p^f=q^{a(\sigma_v)}\quad\text{in }\Z_\ell^\times.
\end{equation}
Its reduction modulo $\ell$ says that $p^f$ generates
$\mathbb F_\ell^\times$. In particular $\gcd(f,\ell-1)=1$.
But $1\leq f\leq[k:\Q]$, and every prime divisor of any $f>1$
in this range divides $m_0$, hence divides $\ell-1$.
Consequently $f=1$.

Taking logarithms in \eqref{eq:generating-norm} gives
\[
 (\log\chi_\ell\circ\beta)(\sigma_v)
   =\log p
   =\log\chi_{k,\ell}(\sigma_v).
\]
The difference is an element of $\X{\ell}{k}$ vanishing at
infinitely many primes of absolute residue degree one.
Proposition~\ref{prop:character-detection}(i) proves
\eqref{eq:cyclotomic-logarithms}.
\end{proof}
